\section*{Capítulo \ref{ch:b2:plant-meristems} --- Desenvolvimento vegetativo das plantas: meristemas e crescimento}

\begin{solution}{exo:b2:plant-meristems:1}
Um \omterm{def:b2:plant-meristems:meristem}{meristema} é uma população persistente de células indiferenciadas em
divisão, da qual derivam os tecidos da planta. \omterm{def:b2:plant-meristems:meristem}{Meristema apical
caulinar}: folhas, caule, \omterm{def:b2:angiosperm-reproduction:flower}{flores}; \omterm{def:b2:plant-meristems:meristem}{meristema apical radicular}: a raiz;
\omterm{def:b2:plant-meristems:wood}{câmbio vascular}: xilema secundário (\omterm{def:b2:plant-meristems:wood}{madeira}) e floema; felogênio: a
\omterm{def:b2:plant-meristems:wood}{periderme} (súber) da \omterm{def:b2:plant-meristems:wood}{casca}.
\end{solution}

\begin{solution}{exo:b2:plant-meristems:2}
Coifa (proteção, lubrificação, percepção da gravidade, células que se
desprendem); \omterm{def:b2:plant-meristems:meristem}{meristema} com o \omterm{prop:b2:plant-meristems:ram}{centro quiescente} (divisão); \omterm{prop:b2:plant-meristems:ram}{zona de
alongamento} (as células se alongam de dez a vinte vezes); zona de
\omterm{def:b2:cell-differentiation:differentiation}{diferenciação} (\omterm{prop:b2:plant-meristems:ram}{pelos radiculares}, xilema e floema maduros). As raízes
laterais surgem do periciclo da raiz madura, em frente aos polos do
xilema.
\end{solution}

\begin{solution}{exo:b2:plant-meristems:3}
Plantas de feijão foram decapitadas: as gemas laterais brotaram. Num
segundo grupo, o coto cortado recebeu um bloco de ágar contendo
auxina: as gemas permaneceram dormentes; no grupo controle, um bloco
de ágar puro: as gemas brotaram. Conclusão: é a auxina vinda do ápice
que inibe as gemas laterais --- a \omterm{prop:b2:plant-meristems:apical}{dominância apical} é hormonal.
\end{solution}

\begin{solution}{exo:b2:plant-meristems:4}
Sessenta anos; o alburno corresponde aos últimos quinze anéis; estão
vivos o \omterm{def:b2:plant-meristems:meristem}{câmbio}, o floema, as células dos raios e do parênquima do
alburno e o felogênio --- uma fina camada em torno de um núcleo morto.
\end{solution}

\begin{solution}{exo:b2:plant-meristems:5}
$0.1\times(P - 0.3)$: \qty{0.01}{h^{-1}} (duplicação em $\ln 2/0.01 =
\qty{69}{h}$), \qty{0.03}{h^{-1}} (\qty{23}{h}),
\qty{0.05}{h^{-1}} (\qty{14}{h}). A \qty{0.25}{MPa}, a turgescência
fica abaixo do limiar de escoamento: o crescimento para, antes de
qualquer murcha.
\end{solution}

\begin{solution}{exo:b2:plant-meristems:6}
Folhas 1--8 a 0, 137.5, 275, 52.5, 190, 327.5, 105, \qty{242.5}{\degree}.
Em ordem: 0, 52.5, 105, 137.5, 190, 242.5, 275, 327.5 --- a menor
lacuna é de \qty{32.5}{\degree}. Cada folha fica numa lacuna das
anteriores, de modo que as oito se sombreiam o mínimo possível e a copa
preenche uniformemente o disco de luz.
\end{solution}

\begin{solution}{exo:b2:plant-meristems:7}
$2\pi r h\Delta r$: com \qty{5}{cm}, $2\pi\times 0.05\times 12\times
0.003 = \qty{0.011}{m^3}$, \qty{5.7}{kg}, \qty{2.8}{kg} de carbono; com
\qty{40}{cm}, \qty{0.090}{m^3}, \qty{45}{kg}, \qty{23}{kg} de carbono
--- oito vezes mais com a mesma largura de anel.
\end{solution}

\begin{solution}{exo:b2:plant-meristems:8}
\emph{clv3}: sem freio --- \omterm{def:b2:plant-meristems:meristem}{meristemas} enormes, folhas e órgãos florais
extras, caules fasciados. \emph{wus}: sem acelerador --- o \omterm{def:b2:plant-meristems:meristem}{meristema} se
esgota após algumas folhas e é refundado repetidamente. Duplo mutante:
o \omterm{def:b2:meiosis-heredity:vocabulary}{fenótipo} \emph{wus}, de modo que o WUS age a jusante do CLV3: a função
do CLV3 é reprimir o WUS, e sem WUS não há nada a reprimir.
\end{solution}

\begin{solution}{exo:b2:plant-meristems:9}
Transporte: \qty{20}{h}. Difusão: $(0.2)^{2}/(2\times 10^{-9}) =
\qty{2e7}{s}$, cerca de 230 dias. O transporte polar torna um sinal
hormonal utilizável ao longo do comprimento de uma planta em um dia, e
lhe dá uma direção.
\end{solution}

\begin{solution}{exo:b2:plant-meristems:10}
\qty{3}{cm} por dia de tecido maduro com células de \qty{200}{\micro m}:
\num{150} células por fileira por dia. Cem células em divisão precisam,
portanto, dividir-se cada uma 1.5 vez por dia: um ciclo de cerca de
\qty{16}{h}.
\end{solution}

\begin{solution}{exo:b2:plant-meristems:11}
A giberelina alonga os entrenós; uma planta que não consegue responder
a ela tem entrenós curtos e um caule baixo. Uma adubação nitrogenada
pesada faz um trigo alto crescer ainda mais, com espigas mais pesadas,
até tombar (acamar) e apodrecer; um caule curto e rígido fica de pé e
põe o carbono extra no grão, e não na palha, de modo que o índice de
colheita aumenta. Na natureza, uma planta baixa é ultrapassada e
sombreada pelas vizinhas e perde.
\end{solution}

\begin{solution}{exo:b2:plant-meristems:12}
Os \omterm{def:b2:plant-meristems:meristem}{meristemas} permitem que uma planta acrescente órgãos onde e quando
as condições os favorecem --- folhas em direção à luz, raízes em
direção à água --- e abandone os que não compensam; toda gema axilar é
um ponto de crescimento de reserva, de modo que ela sobrevive ao
pastejo, ao fogo e à poda rebrotando; e, como os \omterm{def:b2:plant-meristems:meristem}{meristemas} são
perpetuamente embrionários, uma \omterm{def:b2:asexual-reproduction:clone}{geneta} não tem duração de vida fixa. O
custo: uma planta não pode se mover nem reorganizar o que construiu ---
um órgão, uma vez colocado, fica onde está, e um corpo construído por
acréscimo precisa carregar consigo seu passado morto (o cerne).
\end{solution}

\begin{solution}{pb:b2:plant-meristems:1}
\textbf{1.} $150/3 = 50$ folhas por ápice; \num{10000} para a árvore.
\textbf{2.} 0, 137.5, 275, 52.5 e \qty{190}{\degree}.
\textbf{3.} $8\times 137.5 = 1100 = 3\times 360 + 20$; $13\times 137.5
= 1787.5 = 5\times 360 - 12.5$; $21\times 137.5 = 2887.5 = 8\times 360
+ 7.5$: a 22.ª folha é a primeira a menos de \qty{10}{\degree} da
primeira.
\textbf{4.} Cada folha nova fica na lacuna mais larga, de modo que as
folhas se sombreiam o mínimo; é a auxina, cujo transporte polar
(transportadores PIN) drena os arredores de cada primórdio, que fixa o
espaçamento.
\textbf{5.} $\num{10000}\times 80\,\text{cm}^{2} = \qty{80}{m^2}$: a
copa inteira é formada pelas folhas desta estação.
\textbf{6.} Toda ela; as primeiras folhas são construídas com o amido
armazenado na \omterm{def:b2:plant-meristems:wood}{madeira} e nas raízes no ano anterior.
\textbf{7.} Um ápice triplicado forma primórdios mais depressa --- dois
ou três por plastocrono ---, de modo que as folhas e a área foliar
poderiam triplicar, mas os caules seriam fasciados, com hastes em
forma de fita e filotaxia desordenada, e boa parte da folha extra se
sombrearia.
\textbf{8.} Dia: $0.05\times(0.55 - 0.35) = \qty{0.01}{h^{-1}}$; noite:
$0.05\times 0.40 = \qty{0.02}{h^{-1}}$.
\textbf{9.} Dia: $2e^{0.12} = \qty{2.25}{cm}$, \qty{0.25}{cm}
acrescentados; noite: $2.25e^{0.24} = \qty{2.87}{cm}$, \qty{0.61}{cm}
acrescentados.
\textbf{10.} Taxa média de \qty{0.015}{h^{-1}} em \qty{96}{h}: $2e^{1.44}
= \qty{8.4}{cm}$.
\textbf{11.} Dia: $0.05\times 0.30 = \qty{0.015}{h^{-1}}$; noite:
$0.05\times 0.50 = \qty{0.025}{h^{-1}}$.
\textbf{12.} $0.05\times 0.05 = \qty{0.0025}{h^{-1}}$, um quarto da taxa
diurna normal. O ajuste osmótico eleva o conteúdo de solutos da célula
de modo que, ao mesmo potencial hídrico, a turgescência volte a
ultrapassar o limiar.
\textbf{13.} De dia, a transpiração reduz o potencial hídrico da folha
e, com ele, a turgescência, de modo que $P - Y$ é pequeno; à noite, a
turgescência se restabelece. O crescimento é limitado pela água, e não
pelo açúcar, hora a hora.
\textbf{14.} $2\pi\times 0.10\times 8\times 0.005 = \qty{0.025}{m^3}$;
\qty{12.6}{kg} de matéria seca.
\textbf{15.} \qty{6.3}{kg} de carbono.
\textbf{16.} $80\times 4\times 150 = \qty{48}{kg}$ de carbono.
\textbf{17.} $6.3/48 = \qty{13}{\%}$. Outros drenos: as próprias folhas,
as raízes, os galhos e os ramos, a respiração de todas as células vivas
(cerca de metade do que é fixado), as reservas, as \omterm{def:b2:angiosperm-reproduction:flower}{flores} e as
\omterm{def:b2:angiosperm-reproduction:seed}{sementes}.
\textbf{18.} Raio de \qty{0.20}{m}: $2\pi\times 0.20\times 8\times
0.005 = \qty{0.050}{m^3}$, o dobro deste ano --- a circunferência ao
longo da qual o \omterm{def:b2:plant-meristems:meristem}{câmbio} trabalha dobrou (e a árvore estará mais alta).
\textbf{19.} Um ano seco ou frio; ou uma desfolha por insetos, geada ou
fogo. Uma causa climática estreita os mesmos anéis em todas as árvores
da região e se correlaciona com os registros meteorológicos; um dano
afeta uma árvore ou um povoamento, e a geada ou o fogo deixam cicatrizes
anatômicas no anel.
\textbf{20.} As gemas mais próximas de cada corte brotam em poucos dias
e formam vários ramos; em um mês, a árvore está mais ramificada e
densa, com mais ramos, porém mais curtos.
\textbf{21.} A auxina no corte substitui o sinal do ápice: as gemas
permanecem dormentes e não há ramificação.
\textbf{22.} Cada poda remove os ápices e libera as gemas abaixo deles;
seis vezes por ano, cada liberação acrescenta um nível de ramos curtos,
e a superfície se preenche com eles.
\textbf{23.} De gemas axilares dormentes no toco e de gemas
adventícias formadas pelo \omterm{def:b2:plant-meristems:meristem}{câmbio}; cada gema é um \omterm{def:b2:plant-meristems:meristem}{meristema} completo,
capaz de reconstruir o sistema caulinar, ao passo que a cabeça de um
animal é um \omterm{def:b2:vertebrate-development:induction}{organizador} único e insubstituível, sem cópias de reserva.
\textbf{24.} Mais citocinina chegando às gemas promove sua brotação ---
regar uma planta podada a faz ramificar mais.
\textbf{25.} 50 folhas por ápice, \num{10000} por árvore; taxas do
entrenó de \qty{0.01}{h^{-1}} de dia e \qty{0.02}{h^{-1}} à noite; a
\omterm{def:b2:plant-meristems:wood}{madeira} do ano, \qty{0.025}{m^3}, \qty{12.6}{kg} de matéria seca,
\qty{6.3}{kg} de carbono.
\end{solution}
